
% ════════════════════════════════════════════
% 具体分析（5.1.1.具体分析.tex，三级子块）—— 写作提示词见本文件末尾【本章节写作提示词】。
% ════════════════════════════════════════════
\subsubsection{具体分析}

针对问题一，本问属于空间几何可行性判定与单目标连续优化类问题，采用三重约束几何建模与网格数值验证相结合的方法进行求解。

首先对题面给出的建筑、车辆与告警参数进行结构化与坐标换算，接着利用圆盘模型与长方体棱柱模型分别刻画车辆几何约束与视线遮挡约束，然后构建以水平距离最小化为目标的最优观测站模型并解析俯仰约束给出的距离下界，最后通过 0.5 m 网格数值搜索与圆弧细化得到最优站及并列取舍规则。

具体分析流程如图\ref{fig:analysis1_flow}所示。

\begin{figure}[ht]
    \begin{center}
        \begin{tikzpicture}[x=1cm, y=0.7cm]
            \node[box] (n11) at (0, 0) {读入告警参数};
            \node[box] (n12) at (5, 0) {建立三重约束};
            \node[box] (n13) at (10, 0) {网格离散区域};
            \node[box] (n22) at (5, -3) {俯仰距离下界};
            \node[box] (n21) at (0, -3) {可观测性判定};
            \node[box] (n23) at (10, -3) {遮挡判交};
            \node[box] (n31) at (0, -6) {求最小距离};
            \node[box] (n32) at (5, -6) {并列取舍};
            \node[box] (n33) at (10, -6) {输出最优站};
            \draw[arrow] (n11) -- (n12);
            \draw[arrow] (n12) -- (n13);
            \draw[arrow] (n23) -- (n22);
            \draw[arrow] (n22) -- (n21);
            \draw[arrow] (n13) -- ++(0,-1.0) -| (n23);
            \draw[arrow] (n21) -- (n31);
            \draw[arrow] (n31) -- (n32);
            \draw[arrow] (n32) -- (n33);
        \end{tikzpicture}
        \caption{问题一分析流程图}
        \label{fig:analysis1_flow}
    \end{center}
\end{figure}
